% Generated draft: requires mathematical and source review.
\documentclass[11pt]{article}
\usepackage[margin=25mm]{geometry}
\usepackage{fontspec}
\setmainfont{Latin Modern Roman}
\usepackage{amsmath,amssymb,mathtools,unicode-math}
\setmathfont{Latin Modern Math}
\usepackage{enumitem,needspace,xurl,hyperref}
\setlength{\parindent}{0pt}
\setlength{\parskip}{5pt}
\setlength{\emergencystretch}{3em}
\newcommand{\N}{\mathbb{N}}
\newcommand{\Z}{\mathbb{Z}}
\newcommand{\Q}{\mathbb{Q}}
\newcommand{\R}{\mathbb{R}}
\newcommand{\C}{\mathbb{C}}
\newcommand{\F}{\mathbb{F}}
\newcommand{\A}{\mathbb{A}}
\newcommand{\PP}{\mathbb{P}}
\newcommand{\E}{\mathbb{E}}
\newcommand{\eps}{\varepsilon}
\newcommand{\dd}{\,\mathrm{d}}
\begin{document}
\section*{UnsolvedMath ID: 30005550}
\textbf{Source problem code:} catalogue entry
\subsection{Definitions and mathematical statement}
For the standard radial-projection notation, let \(d\geq 2\) be an integer, let \(S^{d-1}=\{\omega\in\R^d:|\omega|=1\}\), and let \(\sigma_{d-1}\) denote surface measure on this sphere. For \(y\in\R^d\), define
\[
\pi_y:\R^d\setminus\{y\}\longrightarrow S^{d-1},\qquad \pi_y(x)=\frac{x-y}{|x-y|}.
\]
For a Borel measure \(\mu\) with \(y\notin\operatorname{supp}\mu\), its radial projection is the pushforward \((\pi_y)_\#\mu\). A finite nonnegative Borel measure \(\lambda\) is called \(u\)-Frostman if there is a constant \(C_\lambda<\infty\) such that
\[
\lambda(B(x,r))\leq C_\lambda r^u\qquad(x\in\R^d,\ r>0).
\]
This gives a conventional meaning to Frostman measure; the supplied excerpt does not specify finiteness, normalization, compactness, or the precise permitted exponent domains, so these conventions require confirmation before fixing a definitive theorem.

When \((\pi_y)_\#\mu\) is absolutely continuous with respect to \(\sigma_{d-1}\), write
\[
f_y=\frac{\dd((\pi_y)_\#\mu)}{\dd\sigma_{d-1}}.
\]
For finite exponents \(1\leq p,q<\infty\), the displayed mixed estimate has the meaning
\[
\int_{\R^d}\|f_y\|_{L^p(S^{d-1},\sigma_{d-1})}^{q}\,\dd\nu(y)<\infty.
\]
Its interpretation includes existence of the indicated density for \(\nu\)-almost every \(y\). The restriction to finite exponents here explains the displayed integral; it is not a supplied restriction on the full research target.

\textbf{Statement to be made precise.} Determine the optimal exponent range for mixed radial-projection estimates of this form for Frostman measures, and identify geometric applications in which improving that range strengthens the conclusion. The supplied excerpt announces a theorem attributed there to reference [7], with hypotheses that \(\mu\) and \(\nu\) are respectively \(s\)- and \(t\)-Frostman measures in \(\R^d\),
\[
\operatorname{supp}\mu\cap\operatorname{supp}\nu=\varnothing,
\qquad s+t>2(d-1),
\]
and with the displayed mixed-integrability conclusion. However, the excerpt ends before giving the restrictions relating \(p,q,s,t,d\). It therefore supplies neither a complete theorem nor enough information to specify the proposed optimal range, its endpoints, or a uniform quantitative bound. The displayed dimension inequality is a hypothesis of that announced theorem; the record does not establish that it is the entire domain of the open question. Disjoint supports alone must not be replaced silently by compact supports or positive separation.

\textbf{Supplied motivation and qualifications.} The excerpt recalls that for a Borel set \(E\subset\R^2\) with \(\dim_H E>1\), the orthogonal projection \(\pi_e(x)=x\cdot e\) has an image of positive one-dimensional Lebesgue measure for almost every \(e\in S^1\). It attributes to Kaufman's 1968 argument the existence of a Frostman measure on \(E\) whose orthogonal projections have \(L^2\) densities for almost every direction. It mentions analogous higher-dimensional results for projections \(\pi_V:\R^d\to V\), where \(V\in G(d,n)\), the Grassmannian of \(n\)-dimensional linear subspaces. Having an \(L^2\) density is sufficient, but is not necessary, for positivity of the projected set's measure. These are supplied background claims, not assertions of independent verification. The excerpt's formula for the range attributed to Orponen in reference [4] is corrupted; only its textual qualification that this range is always below \(2\) can be retained without reconstructing missing mathematics.
\subsection{Short English statement}
The problem asks which mixed integrability exponents are optimal for radial projections of Frostman measures, and where better exponents yield stronger geometric results. The supplied theorem fragment assumes disjoint supports and \(s+t>2(d-1)\), but omits the admissible \(p,q\) range and other details needed for a definitive formulation. The supplied status is partially solved; this statement-only draft requires mathematical and source review.
\subsection{Sources}
\begin{itemize}
\item Local UnsolvedMath registry: supplied record with permanent ID \(30005550\), title \(L^p\) Ranges for Radial Projection Estimates, exact source problem code \texttt{catalogue entry}; collection: Oberwolfach Reports Open Problems; proposed year: 2023. No individual proposer is identified.
\item Supplied report citation: \emph{Incidence Problems in Harmonic Analysis, Geometric Measure Theory, and Ergodic Theory} (2023), Oberwolfach Reports, pp. 1397--1452, DOI \texttt{10.4171/owr/2023/25}. \url{https://doi.org/10.4171/owr/2023/25}. The extracted text includes a page marker for page 1440.
\item Supplied literature source, labelled ``OWR 2023 25'': \url{https://publications.mfo.de/bitstream/handle/mfo/4055/OWR_2023_25.pdf?sequence=1}.
\item Supplied literature source, labelled ``Dąbrowski Orponen Villa'': \url{https://arxiv.org/abs/2107.04471}. No title or explicit correspondence with the excerpt's numbered references is supplied.
\item Literature assessment dated exactly \texttt{2026-08-22}: the supplied curation labels the problem partially solved and states that cited projection and mixed norm theorems establish important \(p\) ranges, while the proposed general improvement of radial projection estimates remains open. This is supplied dated curation, not an independently verified literature assessment.
\item The registry labels the statement \texttt{corrected\_verified} and reports checking it against the official report. This records supplied provenance only; the sources and that verification claim have not been independently checked here. The excerpt's references [4], [5], [6], and [7] lack their bibliographic entries in the supplied record.
\end{itemize}
\end{document}
